% Eighth Philippic --- headnote
% philippic-8, 3 February 43 BC, Rome

\textit{Delivered in the Senate on 3 February 43 BC, the morning
after the legation despatched in early January had returned from
Antony at Mutina, the Eighth Philippic is Cicero's answer to a
vote he had just lost. On 2 February the Senate, presided over by
the consul Vibius Pansa, had debated how to name the crisis. Cicero
and the firmer party wanted the word \textit{bellum} --- war ---
set down in the decree; Lucius Caesar, Antony's uncle, moved the
softer line, and the Senate decreed a \textit{tumultus} (a state
of emergency) and the putting-on of the war-cloaks, but balked at
the word ``war.'' Philippic 8 reopens the question the next day,
and it opens with a rare thing in the series: an open rebuke of the
presiding consul. ``Yesterday's business was conducted more
confusedly, Gaius Pansa, than the practice of your consulship
demanded'' --- Pansa let the men he is not used to yielding to
carry the milder vote.}

\textit{The speech turns first on a point of language
(\textsection 1--4). Cicero will not let the euphemism stand,
and he draws the famous distinction: there can be a war without a
tumult, but no tumult without a war; a \textit{tumultus} is the
graver word, not the lesser, since in war the legal exemptions
from service hold good and in a tumult they do not. The verbal
quibble has a hard practical edge: take away the name of war and
you take away the zeal of the municipalities and the consensus of
the Roman people, which has already committed itself to the cause.
He then proves the thing is a war whatever it is called
(\textsection 5--8): Decimus Brutus is besieged in Mutina, Gaul
is laid waste, Hirtius has already dislodged Antony's garrison
from Claterna and Octavian has taken the field of his own accord.
This is, Cicero reckons, the fifth civil war of his lifetime, but
the first fought not in the discord of citizens but in an
incredible concord --- the whole state, for its temples, walls,
hearths, and the tombs of its ancestors, against one man. The
antithesis of promises follows (\textsection 9--10): what Antony
holds out to his bandits (the city divided up among them, the
plunder of the great houses, Caesar's auction-spear) set against
what the Senate promises its soldiers (liberty, laws, courts, the
empire of the world, peace).}

\textit{The long centre of the speech is a duel, conducted by
name, with Quintus Fufius Calenus (\textsection 11--19), the
consular who had praised the advantages of peace and made himself
Antony's most steadfast defender. Cicero's question is the hinge:
do you call slavery peace? He marshals the canonical exempla of
justified senatorial force --- Scipio Nasica against Tiberius
Gracchus, Opimius armed under the first ``last decree'' against
Gaius Gracchus, Marius and the \textit{optimates} against
Saturninus, his own consulship against Catiline --- and answers
Calenus's principle that all citizens should be kept safe with the
surgeon's metaphor: as a diseased limb is burned and cut away to
save the body, so what is pestilent in the body of the
commonwealth must be amputated. The exchange has its famous barb:
Cicero ``concedes'' that in one man, Publius Clodius, Calenus saw
more truly than he did --- a sarcasm built from a string of
inverted virtues. Then the indictment of the embassy itself
(\textsection 20--28): the shame of sending legates a second time
while Antony battered Mutina before their very eyes, set against
the model of Gaius Popilius Laenas, who drew a circle in the sand
around King Antiochus and would not let him step out of it until
he had answered the Senate. Cicero reads out Antony's ``modest''
counter-demands --- the provinces, lands for his six legions, the
untouched accounts at the temple of Ops, protection for his
creatures Cafo, Saxa, Nucula, Mustela, and Tiro, and his judiciary
law --- and contrasts the contempt Antony showed Rome's legates
with the deference Rome paid to his legate Cotyla.}

\textit{The close turns to the duty of a leading man
(\textsection 29--32): a princeps must serve the citizens not
only in their minds but in their very eyes, and Cicero offers the
aged augur Quintus Scaevola, who in the Marsic War, old and
broken in health, was the first into the Senate-house every
morning, as the model the timid consulars should imitate or at
least not envy. He announces that he will not claim the consular's
privilege of keeping the toga while the city is in war-cloaks, but
will dress as every other citizen does. The speech ends, as the
deliberative speeches do, in a formal motion (\textsection 33):
amnesty for any of Antony's men who lay down their arms and come
over to the consuls, to Decimus Brutus, or to Octavian before the
next Ides of March; rewards to be referred to the Senate; and the
brand of public enemy on anyone who goes over to Antony after the
decree --- with the single bitter exception of Varius (Cotyla),
whom Cicero lets return to his commander on the one condition that
he never set foot in Rome again. Where Philippic 7 held the line
during the legation interval, Philippic 8 is the speech of the
morning after: the legation Cicero had opposed has come back with
nothing but Antony's insolence, exactly as he foretold, and he
will not let the Senate hide the war behind a milder word.}
